\section{Spreading Fourier coefficients}
\label{sec:spreading}

This section supplies the auxiliary polynomial used in the sampling
construction.  We reproduce the coefficient-spreading argument of
\cite[Section~4]{OpenAIMinimal}, constructing a real trigonometric
polynomial with almost unit mean-square mass whose Fourier coefficients
fit within a total spectral width less than one.
The coefficient assigned to a width $w$ must have size at
most approximately $\sqrt w$: this is the normalization of the quadratic
integral in Lemma~\ref{lem:quadratic}.  Additional torus variables let us
distribute each coefficient among many frequencies while controlling
both their number and their sizes.

For a trigonometric polynomial $F$ on $\T^m$, write
\[
 \widehat F(a)=\int_{\T^m}F(x)\e(-a\cdot x)\,dx,
 \qquad a\in\Z^m.
\]
A finite set is a \emph{containing Fourier support} if every coefficient
outside it vanishes; coefficients inside it are allowed to vanish as well.
All norms and integrals on tori use probability Haar measure.

\begin{proposition}[An auxiliary polynomial with controlled widths]
\label{prop:auxiliary}
For every $0<\delta<1/20$, there exist integers $m\ge1$ and $r\ge2$,
nonzero vectors $a_1,\ldots,a_r\in\Z^m$, a real-valued trigonometric
polynomial $F$ on $\T^m$, and a vector $v\in\R^m$ with the following
properties.  No two of the vectors $a_i$ are parallel, and
\[
 \mathcal A=\{a_1,-a_1,\ldots,a_r,-a_r\}
\]
is a containing Fourier support of $F$.  Moreover,
\begin{equation}
 \label{eq:auxiliary-norms}
 \norm{F}_\infty\le1,
 \qquad \norm{F}_2\ge1-3\delta.
\end{equation}
Writing $\lambda_a=a\cdot v$ for $a\in\mathcal A$, we have
\begin{equation}
 \label{eq:auxiliary-widths}
 \lambda_a\ne0,
 \qquad \sum_{a\in\mathcal A}|\lambda_a|<1,
 \qquad
 \frac{|\widehat F(a)|}{\sqrt{|\lambda_a|}}
 \le K_\delta,
 \qquad
 K_\delta=\sqrt{\frac{(1+\delta)^3}{1-\delta}}.
\end{equation}
All these objects depend on $\delta$ alone.
\end{proposition}

The construction starts with a polynomial of almost unit mean square.
Quadratic oscillations in new variables spread its coefficients; finite
Fourier truncation then gives the required polynomial.  We first record
the uniform quadratic estimate used to control the spread coefficients.

\begin{lemma}[Uniform quadratic integration]
\label{lem:quadratic}
Let $g\in C_c^\infty(\R)$ and let $\beta\in\R\setminus\{0\}$ be fixed.
As $T\to\infty$ through positive real values,
\begin{equation}
 \label{eq:uniform-quadratic}
 \begin{split}
 \int_{\R}g(x)\e(T\beta x^2/2-ux)\,dx
 ={}&\frac{\e(\operatorname{sgn}(\beta)/8-u^2/(2T\beta))}
 {\sqrt{T|\beta|}}\\
 &\qquad\cdot\left(g\left(\frac{u}{T\beta}\right)
       +O_{g,\beta}(T^{-1})\right),
 \end{split}
\end{equation}
uniformly for every real $u$.
\end{lemma}

\begin{proof}
Use the real-line Fourier transform
$\widehat g(\xi)=\int_{\R}g(x)\e(-\xi x)\,dx$, and put
$x_0=u/(T\beta)$.  Completing the square reduces the integral to
\[
 \e\left(-\frac{u^2}{2T\beta}\right)
 \int_{\R}g(x_0+y)\e(T\beta y^2/2)\,dy.
\]
Insert the damping factor $\exp(-\pi\rho y^2)$ with $\rho>0$, and use
Fourier inversion on $g(x_0+y)$.  The resulting double integral is
absolutely integrable.  The inner Gaussian integral is
\[
 \int_{\R}\exp\bigl(-\pi(\rho-iT\beta)y^2+2\pi i\xi y\bigr)\,dy
 =(\rho-iT\beta)^{-1/2}
   \exp\left(-\frac{\pi\xi^2}{\rho-iT\beta}\right).
\]
Here the square root is the branch analytic in the right half-plane
and positive on the positive real axis.  As $\rho\downarrow0$,
\[
 (\rho-iT\beta)^{-1/2}
 \longrightarrow\frac{\e(\operatorname{sgn}(\beta)/8)}{\sqrt{T|\beta|}},
 \qquad
 \exp\left(-\frac{\pi\xi^2}{\rho-iT\beta}\right)
 \longrightarrow\e\left(-\frac{\xi^2}{2T\beta}\right).
\]
The modulus of the exponential multiplier is at most one, and
$|(\rho-iT\beta)^{-1/2}|\le(T|\beta|)^{-1/2}$.  Since $\widehat g$ is
Schwartz, dominated convergence applies on the Fourier side.  On the
original side it applies with dominating function $|g(x_0+y)|$.
We obtain the exact identity
\begin{equation}
 \label{eq:quadratic-exact}
 \begin{split}
 \int_{\R}g(x)\e(T\beta x^2/2-ux)\,dx
 ={}&\frac{\e(\operatorname{sgn}(\beta)/8-u^2/(2T\beta))}
 {\sqrt{T|\beta|}}\\
 &\qquad\cdot\int_{\R}\widehat g(\xi)\e(\xi x_0)
                   \e\left(-\frac{\xi^2}{2T\beta}\right)d\xi.
 \end{split}
\end{equation}
The difference of the last integral from $g(x_0)$ has absolute value at
most
\[
 \frac{\pi}{T|\beta|}\int_{\R}\xi^2|\widehat g(\xi)|\,d\xi.
\]
This bound is independent of $x_0$, proving the required uniformity.
\end{proof}

\begin{proof}[Proof of Proposition~\ref{prop:auxiliary}]
Fix $\delta$.  The choices below are successive: once a parameter has
been chosen, it stays fixed while later parameters vary.

\paragraph{A polynomial of almost unit mean square.}
Starting with $p_0=0$, introduce one new torus coordinate at each step
and define
\begin{equation}
 \label{eq:binary-recursion}
 p_j(y_1,\ldots,y_j)
 =p_{j-1}(y_1,\ldots,y_{j-1})
  +\frac{1-p_{j-1}(y_1,\ldots,y_{j-1})^2}{2}\cos(2\pi y_j).
\end{equation}
For $|z|\le1$,
$(1-z^2)/2=(1-|z|)(1+|z|)/2\le1-|z|$.  Thus $|p_j|\le1$
pointwise.  Integrating first in $y_j$ shows that $p_j$ has mean zero.
If $M_j=\int_{\T^j}p_j^2$, integration of the square in the new variable
gives
\begin{equation}
 \label{eq:binary-energy}
 M_j=M_{j-1}+\frac18\int_{\T^{j-1}}(1-p_{j-1}^2)^2
 \ge M_{j-1}+\frac18(1-M_{j-1})^2.
\end{equation}
Hence $M_j\uparrow1$: a limit below one would make the increments
bounded below by a positive constant.  Choose $d\ge2$ so that
$M_d\ge1-\delta$.

Let $\mathcal S\subset\Z^d$ be the Fourier support of $p_d$, consisting
of the frequencies with nonzero coefficients, and put
$c_s=\widehat p_d(s)$.  Every vector in $\mathcal S$ has last nonzero
coordinate equal to $1$ or $-1$.  Indeed, the old term in
\eqref{eq:binary-recursion} keeps the previous frequencies with new
coordinate zero, whereas every frequency in the new term has new
coordinate $1$ or $-1$.  The two terms cannot cancel each other.
There is no constant frequency because the mean is zero.

If two such vectors are parallel, comparing their last nonzero
coordinates shows that they are equal or negatives.  Let
$s_1,\ldots,s_D$ be those with positive last nonzero coordinate.
They are pairwise nonparallel and
$\mathcal S=\{\pm s_1,\ldots,\pm s_D\}$.  There are at least two:
the coefficient at the first coordinate vector in $p_1$ is nonzero,
and the coefficient at the second coordinate vector in $p_2$ is
$(1-M_1)/4=7/32>0$; both persist in later steps.

We have obtained the mean-square mass needed in
\eqref{eq:auxiliary-norms}.  It remains to distribute its Fourier
coefficients so that the total width and the coefficient-to-width
ratios in \eqref{eq:auxiliary-widths} are both controlled.

\paragraph{Choosing the quadratic curvatures.}
Choose $\gamma\in\R^d$ outside the finite union of hyperplanes
$s^\perp$, $s\in\mathcal S$, and define
$\rho_s=|s\cdot\gamma|>0$.  We will choose vectors
$b_1,\ldots,b_D\in\R^d$ so that the products
\[
 \Delta_s=\prod_{j=1}^D|s\cdot b_j|
\]
are approximately a common multiple of $|c_s|^2/\rho_s$.
Here is the reason for this target. At a large oscillation scale $B$,
the construction below assigns about $B^D\Delta_s$ Fourier positions
to the original frequency $s$, with coefficient sizes bounded by about
$|c_s|/\sqrt{B^D\Delta_s}$. Assigning each such position a width
proportional to $\rho_s/B^D$ makes the total width of this family
proportional to $\Delta_s\rho_s$, and the resulting upper bound for the
ratio of squared coefficient size to width proportional to
$|c_s|^2/(\Delta_s\rho_s)$.
Thus choosing $\Delta_s\rho_s$ approximately proportional to $|c_s|^2$
balances these upper bounds while making the total width track the Parseval mass
$\sum_s|c_s|^2$. The estimates and the common proportionality factor
are made precise below.

For each $j$, choose $b_j^0\in s_j^\perp$ which is not orthogonal to
any $s_l$ with $l\ne j$.  This is possible because each such
orthogonality condition cuts out a proper subspace of $s_j^\perp$,
and a finite union of proper subspaces cannot cover that space.
Choose $b_j^1$ satisfying
\[
 s_j\cdot b_j^1
 =\frac{|c_{s_j}|^2}
 {\rho_{s_j}\prod_{l\ne j}|s_j\cdot b_l^0|}>0,
\]
and set $b_j=b_j^0+\tau b_j^1$ for a positive parameter $\tau$.
Every denominator is nonzero by the choices just made.  For $s=s_j$,
\[
 \frac{\Delta_{s_j}}\tau
 =(s_j\cdot b_j^1)
   \prod_{l\ne j}|s_j\cdot b_l^0+\tau s_j\cdot b_l^1|
 \longrightarrow\frac{|c_{s_j}|^2}{\rho_{s_j}}
 \qquad(\tau\downarrow0).
\]
The same limit holds for $-s_j$.  Fix a sufficiently small $\tau>0$
so that, simultaneously for all $s\in\mathcal S$,
\begin{equation}
 \label{eq:determinant-tuning}
 (1-\delta)\frac{\tau|c_s|^2}{\rho_s}
 \le\Delta_s\le
 (1+\delta)\frac{\tau|c_s|^2}{\rho_s}.
\end{equation}
In particular, every curvature
$\beta_{s,j}:=s\cdot b_j$ is nonzero.  All these curvatures are fixed
from now on; no lower bound uniform in $\delta$ is needed.

\paragraph{Spreading the coefficients.}
Choose $g\in C_c^\infty((-1/2,1/2))$ with $0\le g\le1$ and
\[
 \left(\int_{\R}g(x)^2\,dx\right)^D\ge1-\delta.
\]
For a large positive real parameter $B$, and
$z\in[-1/2,1/2]^D$, define
\begin{equation}
 \label{eq:spread-function}
 G(y,z)=\left(\prod_{j=1}^Dg(z_j)\right)
 p_d\left(y+\frac B2\sum_{j=1}^Db_jz_j^2\right),
 \qquad y\in\T^d.
\end{equation}
The cutoff makes this function vanish near every boundary face in
$z$, so it extends to a smooth real function on $\T^{d+D}$.
Translation invariance in $y$ gives
\begin{equation}
 \label{eq:spread-norms}
 \norm{G}_\infty\le1,
 \qquad
 \norm{G}_2^2=M_d\left(\int g^2\right)^D\ge(1-\delta)^2.
\end{equation}

Fourier integration first in $y$ shows that the only possible
frequencies are $(s,u)$ with $s\in\mathcal S$ and $u\in\Z^D$, and
\begin{equation}
 \label{eq:spread-coefficients}
 \widehat G(s,u)=c_s\prod_{j=1}^DI_{s,j}(u_j),
 \qquad
 I_{s,j}(w)=\int_{\R}g(x)\e(B\beta_{s,j}x^2/2-wx)\,dx.
\end{equation}
Lemma~\ref{lem:quadratic} and $0\le g\le1$ imply
\[
 |I_{s,j}(w)|\le
 \frac{1+O(B^{-1})}{\sqrt{B|\beta_{s,j}|}}
\]
uniformly for every real $w$.  The constants may depend on the
already fixed curvatures and cutoff.  Since there are only finitely
many pairs $(s,j)$, for all sufficiently large $B$ we have
\begin{equation}
 \label{eq:spread-coefficient-bound}
 |\widehat G(s,u)|\le
 (1+\delta)\frac{|c_s|}{\sqrt{B^D\Delta_s}}
 \qquad(s\in\mathcal S,\ u\in\Z^D).
\end{equation}

The stationary points for these integrals suggest keeping the boxes
\[
 \mathcal U_s=
 \left\{u\in\Z^D: |u_j|\le\frac{B|\beta_{s,j}|}{2}
       \text{ for every }j\right\}.
\]
Their cardinalities satisfy
\begin{align}
 |\mathcal U_s|
 &=\prod_{j=1}^D\left(2\left\lfloor
                  \frac{B|\beta_{s,j}|}{2}\right\rfloor+1\right)
 \notag\\
 &\le B^D\Delta_s\prod_{j=1}^D
       \left(1+\frac1{B|\beta_{s,j}|}\right)
 \le(1+\delta)B^D\Delta_s
 \label{eq:spread-box-count}
\end{align}
for all sufficiently large $B$.  Thus the product $\Delta_s$ controls
the number of retained frequencies as well as their coefficient bound.

\paragraph{A summable bound outside the boxes.}
We need uniform approximation after truncation, so the additive error
in Lemma~\ref{lem:quadratic} cannot simply be summed over the whole
lattice.  Instead, outside the boxes the phase has no stationary
point near the amplitude support.

Choose $\alpha<1/2$ with $\operatorname{supp}g\subset[-\alpha,\alpha]$.
For a fixed $\beta\ne0$, if $|w|>B|\beta|/2$, then
\begin{equation}
 \label{eq:offbox-derivative}
 |B\beta x-w|\ge |w|-B|\beta|\alpha
 \ge c_{\alpha,\beta}(B+|w|)
 \qquad(x\in\operatorname{supp}g),
\end{equation}
with $c_{\alpha,\beta}>0$.  Indeed, the middle expression is at
least $(1-2\alpha)|w|$, while $B<2|w|/|\beta|$.
Integrate by parts $l$ times using
\[
 \frac{1}{2\pi i(B\beta x-w)}\frac{d}{dx}
   \e(B\beta x^2/2-wx)=\e(B\beta x^2/2-wx).
\]
There are no boundary terms.  Each term after these integrations
contains a bounded derivative of $g$ and a factor bounded by
$C_l B^h(B+|w|)^{-l-h}$ for some $0\le h\le l$.  Hence
\begin{equation}
 \label{eq:offbox-bound}
 \left|\int_{\R}g(x)\e(B\beta x^2/2-wx)\,dx\right|
 \le C_l(B+|w|)^{-l}
 \qquad\left(|w|>\frac{B|\beta|}{2}\right).
\end{equation}
For integer $w$ the sum of this bound outside the one-dimensional
box is $O(B^{1-l})$ when $l>1$.  The full sum over integer $w$ is
$O(B)$: inside there are $O(B)$ terms, each at most $\norm{g}_{L^1(\R)}$,
and the outside sum is bounded by \eqref{eq:offbox-bound}.

The complement of $\mathcal U_s$ is contained in the union of the
$D$ coordinate complements.  Consequently
\begin{align}
 \sum_{u\notin\mathcal U_s}|\widehat G(s,u)|
 &\le |c_s|\sum_{j=1}^D
 \left(\sum_{|u_j|>B|\beta_{s,j}|/2}|I_{s,j}(u_j)|\right)
 \prod_{h\ne j}\left(\sum_{u_h\in\Z}|I_{s,h}(u_h)|\right)
 \notag\\
 &=O(B^{D-l}).
 \label{eq:spread-tail}
\end{align}
Fix an integer $l>D$.  Summing over the finite set $\mathcal S$
shows that the total discarded coefficient mass tends to zero.
The Fourier series is absolutely convergent, so this mass also
bounds the uniform truncation error.

\paragraph{Truncation and the final budgets.}
Fix one $B$ large enough for \eqref{eq:spread-coefficient-bound},
\eqref{eq:spread-box-count}, and a truncation error at most $\delta$
in \eqref{eq:spread-tail}.  Define
\[
 T_B(y,z)=\sum_{s\in\mathcal S}\sum_{u\in\mathcal U_s}
              \widehat G(s,u)\e(s\cdot y+u\cdot z),
 \qquad F=\frac{T_B}{1+\delta}.
\]
The boxes satisfy $\mathcal U_{-s}=\mathcal U_s=-\mathcal U_s$.
Since $G$ is real, the retained coefficients have conjugate symmetry,
and $F$ is real.  By \eqref{eq:spread-norms} and the uniform error,
\[
 \norm{F}_\infty\le1,
 \qquad
 \norm{F}_2\ge\frac{1-2\delta}{1+\delta}\ge1-3\delta.
\]

Set $m=d+D$ and take the containing support
\[
 \mathcal A=\{(s,u):s\in\mathcal S,\ u\in\mathcal U_s\}.
\]
Its vectors are nonzero.  If two are parallel, projection onto the
first $d$ coordinates forces the proportionality factor to be $1$
or $-1$; the full vectors must then be equal or negatives.
Thus one representative from each signed pair gives nonparallel
vectors $a_1,\ldots,a_r$.  Every box contains zero, so $r\ge D\ge2$.
Some of the corresponding Fourier coefficients can vanish; the
containing support and all estimates remain valid in that case.

Finally choose
\[
 v=\left(\frac{\gamma}{\tau B^D(1+\delta)^3},0\right)\in\R^{d+D}.
\]
For $a=(s,u)\in\mathcal A$,
\[
 |\lambda_a|=|a\cdot v|
 =\frac{\rho_s}{\tau B^D(1+\delta)^3}>0.
\]
Using \eqref{eq:spread-box-count} and
\eqref{eq:determinant-tuning}, with both signs already included in
$\mathcal S$, gives
\begin{align*}
 \sum_{a\in\mathcal A}|\lambda_a|
 &=\sum_{s\in\mathcal S}|\mathcal U_s|
     \frac{\rho_s}{\tau B^D(1+\delta)^3}\\
 &\le\frac1{\tau(1+\delta)^2}\sum_{s\in\mathcal S}\Delta_s\rho_s
 \le\frac1{1+\delta}\sum_{s\in\mathcal S}|c_s|^2
 =\frac{M_d}{1+\delta}<1.
\end{align*}
The coefficient bound \eqref{eq:spread-coefficient-bound}, divided
by $1+\delta$, also gives
\[
 \frac{|\widehat F(a)|^2}{|\lambda_a|}
 \le\frac{|c_s|^2\tau(1+\delta)^3}{\Delta_s\rho_s}
 \le\frac{(1+\delta)^3}{1-\delta}=K_\delta^2.
\]
These are all the assertions of the proposition.
\end{proof}

The strict signed width bound means that
$w_i=|\lambda_{a_i}|$ satisfies $2\sum_iw_i<1$, exactly the input
required by Lemma~\ref{lem:signed-packing}.  The polynomial $F$, its
support, and its velocity $v$ are now fixed.  After the signed intervals have been packed, Section~\ref{sec:sampling}
will sample these fixed data along quadratic phases while only the length
$N$ tends to infinity.
